% !TEX root=../manuscript.tex
\section{Proofs for Some Lemmas Used in the Paper}

\subsection{Proof for Lemma~\ref{rule-bijection}}
\label{sec:appendix-lemma1}
\begin{lemma}
There is a bijection $\varphi$ between typing rules and synthesis rules.
\end{lemma}
\begin{proof}

$\varphi$ can be proven to be a bijection by showing that it is both surjective and injective.
\begin{itemize}
 \item The translation function $\varphi$ maps every typing rule to a synthesis rule, and synthesis rules are constructed exclusively in this manner; thus, $\varphi$ is surjective.
 \item To prove that $\varphi$ is injective, we define its inverse function $\psi$: given a synthesis rule in the form presented in Figure~\ref{S-Rule}, $\psi$ removes all substitutions $\sigma_i$, as well as the unification and assignment acquisition, and restores the terms $\overline{x}_0$ to the conclusion, recovering the typing rule as shown in Figure~\ref{T-Rule}. Therefore, $\psi \circ \varphi = \varphi \circ \psi = id$. The existence of $\psi$ demonstrates that $\varphi$ is injective.
\end{itemize}

\end{proof}

\subsection{Proof for Lemma~\ref{lemma:assignment}}
\label{sec:appendix-lemma2}
\begin{lemma}
For any synthesis judgment derived using a synthesis rule of the form in Figure~\ref{S-Rule}, the composed substitution $\sigma_c = \sigma_n \ldots \sigma_0 \sigma$ constitutes an assignment, in the sense that $\sigma_c(\skhzc{\overline{x}_i})$, $\sigma_c(\skhzc{\overline{x}})$, and $\sigma_c(\skhzc{\overline{y}})$ are all ground terms.
\end{lemma}
\begin{proof}
Induction on the structure of the synthesis derivation.
\begin{itemize}
 \item \textbf{Base case:} The base case occurs when the synthesis rule has no subgoals. In this situation, the free variables $\overline{x}_f$ are precisely those appearing in the typing rule's conclusion, i.e., $\mvhzc(\overline{x}_f) = \mvhzc(\overline{x}_0)$. The unification step produces a substitution $\sigma$ such that $\sigma(\overline{x}_0) = \sigma(\overline{x})$, and the assignment $\sigma_0$ obtained in this step covers all variables in $\sigma(\overline{x}_f)$. 
 
 Consequently, $\sigma_0\sigma(\overline{x}_f)$ are ground terms, thus $\sigma_0\sigma(\overline{x}_0) = \sigma_0\sigma(\overline{x})$ are also ground terms, and so is $\sigma_0\sigma(\overline{y})$ since $\overline{y}$ contains only variables that are in $\overline{x}_0$.
 
 \item \textbf{Inductive case:} The inductive case arises when the synthesis rule has one or more subgoals. By the induction hypothesis, for each $i$, the substitution $\sigma_i$ is an assignment that covers all variables in $\sigma_{i-1}\ldots\sigma_0\sigma(\overline{x}_i)$, ensuring that $\sigma_{i}\ldots\sigma_0\sigma(\overline{x}_i)$ are ground terms, and therefore so is $\sigma_c(\overline{x}_i)$ for all $i\ge1$.
 
 We similarly observe that $\sigma_0\sigma$ covers all variables in $\overline{x}_f$ so $\sigma_c$ must also cover $\overline{x}_f$ and thus $\overline{x}_0$ ($\overline{x}_f$ contains all variables in $\overline{x}_0$ besides those already instantiated in $\overline{x}_i$). 
 %
 We finally have $\sigma_c(\overline{x}) = \sigma_n \ldots \sigma_0 \sigma(\overline{x}) = \sigma_n \ldots \sigma_0 \sigma(\overline{x}_0) = \sigma_c(\overline{x}_0)$ is a ground term, and so is $\sigma_c(\overline{y})$, since $\overline{y}$ contains only variables that are in $\bigcup_{i=0}^n\overline{x}_i$.
 \end{itemize}
\end{proof}

\subsection{Proof for Lemma~\ref{lemma:translation}}
\label{sec:appendix-lemma3}
\begin{lemma}
Let $T_1$ be a well-formed type derivation tree with root node labeled $P(\overline{x}_c)$, annotated with typing rule $R$, and let $T_2$ be its corresponding synthesis derivation tree constructed via the translation, whose root node is labeled with $P(\overline{x}) \leadsto \theta$, where $\theta$ is the assignment synthesized by Algorithm~\ref{alg:synthtree}.
If $\overline{x}$ is unifiable with $\overline{x}_c$, then in the application of the synthesis rule that derives $P(\overline{x}) \leadsto \theta$, (1) the unification step, (2) the acquisition for free variables, and (3) the constraint checking step will always succeed, and (4) the synthesized assignment $\theta$ makes $P(\overline{x})$ concretized to $P(\theta(\overline{x})) = P(\overline{x}_c)$.
\end{lemma}

\begin{proof}
  Induction on the size of the type derivation tree $T_1$. Suppose the typing rule $R$ has the form shown in Figure~\ref{T-Rule}, and its corresponding synthesis rule $R'$ has the form shown in Figure~\ref{S-Rule}.
  \begin{itemize}
    \item \textbf{Base case:} 
    The base case is when the type derivation tree $T_1$ has only one node, which means the typing rule $R$ has no premises. 
    In this case, since $\overline{x}$ is unifiable with $\overline{x}_c$ and in the typing rule $R$, $\overline{x}_c$ is the concretization of $\overline{x}_0$, the unification and acquisition steps will succeed, yielding a substitution $\sigma$ such that $\sigma(\overline{x}_0) = \sigma(\overline{x})$ and $\theta'\sigma(\overline{x}_0) = \overline{x}_c$ for some assignment $\theta'$, which can be used to obtain $\sigma_0 = \theta'|_{\mvhzc(\sigma(\overline{x}_f))}= \theta'|_{\mvhzc(\sigma(\overline{x}_0))}$.
    
    By Lemma~\ref{lemma:assignment}, $\sigma_0\sigma$ is an assignment that covers all variables in $\overline{x}$, $\overline{x}_0$ and $\overline{y}$, and $\sigma_0\sigma(\overline{x}_0) = \theta'\sigma(\overline{x}_0) = \overline{x}_c$, so if we set $\theta = \sigma_0\sigma|_{\mvhzc(\overline{x})}$ we have $\theta(\overline{x}) = \sigma_0\sigma(\overline{x})= \sigma_0\sigma(\overline{x}_0) = \overline{x}_c$. The constraint checking step is also satisfied because $\mvhzc(y)\subseteq \mvhzc(\overline{x}_0)$ and $\sigma_0\sigma(\overline{y})$ is the same as the verified concretized constraint in the type derivation tree.
    
    \item \textbf{Inductive case:} If the typing rule $R$ has a series of premises $P_1(\overline{x}_1), \ldots, P_n(\overline{x}_n)$, then the synthesis rule $R'$ has a series of subgoals $P_1(\meta{\sigma_0\sigma(\overline{x}_1)}) \leadsto \sigma_1, \ldots, P_n(\meta{\sigma_{n-1}\ldots\sigma_0\sigma(\overline{x}_n)}) \leadsto \sigma_n$.
    
    As in the base case, the unification and acquisition steps will succeed, yielding a substitution $\sigma$ such that $\sigma(\overline{x}_0) = \sigma(\overline{x})$ and $\theta'\sigma(\overline{x}_0) = \overline{x}_c$ for some assignment $\theta'$, and $\sigma_0 = \theta'|_{\mvhzc(\overline{x}_f)}$.
    
    In the type derivation tree, suppose the $i$-th child node of the root node is labeled with $P_i(\overline{x}_{ic})$, then rule $R$ can be instantiated with some assignment $\theta_c$ such that $\theta_c(\overline{x}_i) = \overline{x}_{ic}$ for all $i$ ($\overline{x}_{0c}$ is just $\overline{x}_{c}$).
    %
    We use the notation $\sigma_i \trianglelefteq \sigma_j$ to denote that $\sigma_j$ factors through $\sigma_i$, meaning that $\sigma_j = \sigma'\sigma_i$ for some substitution $\sigma'$.
    Following the derivation of $\sigma_0$, we know that $\theta'\sigma(\overline{x}_0) = \overline{x}_c$ and $\sigma_0 = \theta'|_{\mvhzc(\sigma(\overline{x}_f))}\trianglelefteq \theta'|_{\mvhzc(\sigma(\overline{x}_0))}$, thus $(\sigma_0\sigma|_{\mvhzc(\overline{x}_0)})\trianglelefteq (\theta'\sigma|_{\mvhzc(\overline{x}_0)}) \trianglelefteq \theta_c$ holds ($\theta_c$ concretizes $\overline{x}_0$) and $\sigma_0\sigma$ ranges over only variables in $\overline{x}_0$ and $\overline{x}$.
    
    Since $\overline{x}_1$ has no common variables with $\overline{x}$, $\sigma_0\sigma(\overline{x}_1)= (\sigma_0\sigma|_{\mvhzc(\overline{x}_0)})(\overline{x}_1)$, and $\theta_c(\overline{x}_1)=\overline{x}_{1c}$, $\sigma_0\sigma(\overline{x}_1)$ is unifiable with $\overline{x}_{1c}$.
    %
    Thus we can use the induction hypothesis to synthesize $P_1(\meta{\sigma_0\sigma(\overline{x}_1)}) \leadsto \sigma_1$, where $\sigma_1\sigma_0\sigma(\overline{x}_1) = \overline{x}_{1c}$.
    
    \emergencystretch=3em
    Then, note that $\sigma_1\sigma_0\sigma(\overline{x}_1) = \theta_c(\overline{x}_1)$, and $(\sigma_1\sigma_0\sigma|_{\mvhzc(\overline{x}_0,\overline{x}_1)}) \allowbreak\trianglelefteq \theta'$ holds, and then $\sigma_1\sigma_0\sigma(\overline{x}_2)=\allowbreak(\sigma_1\sigma_0\sigma|_{\mvhzc(\overline{x}_0,\overline{x}_1)})(\overline{x}_2)$ is unifiable with $\overline{x}_{2c}$, and we can apply the induction hypothesis again and again until all assignments to subgoals are synthesized.

    \emergencystretch=0pt
    Finally, we can collect all the assignments $\sigma_n, \ldots, \sigma_1, \sigma_0$ and compose them to form $\sigma_c = \sigma_n \ldots \sigma_0\sigma$, which is an assignment and covers all variables in $\overline{x}_i$, $\overline{x}$ and $\overline{y}$ by Lemma~\ref{lemma:assignment}. $\sigma_c(\overline{x}_i)=\overline{x}_{ic}=\theta_c(\overline{x}_i)$ for all $i$, so $\theta_c \trianglelefteq \sigma_c$ since $\sigma_c$ also ranges over $\mvhzc(\overline{x})$.
    We have $\sigma_c(\overline{y}) = \theta_c(\overline{y})$, which instantiates the constraint $\phi(\overline{y})$ to be satisfied as in the type derivation tree. Here, $\theta(\overline{x}) = \sigma_c|_{\mvhzc(\overline{x})}$ concretizes the synthesis goal $P(\overline{x})$ to the ground typing judgment $P(\overline{x}_c)$ ($\theta(\overline{x})=\sigma_n \ldots \sigma_0\sigma(\overline{x})=\sigma_n \ldots \sigma_0\sigma(\overline{x}_0)=\sigma_c(\overline{x}_0)=\overline{x}_c$), which is the same as the root node of the type derivation tree.
  \end{itemize}
\end{proof}

\subsection{Proof for Lemma~\ref{lemma:isomorphism2}}
\label{sec:appendix-lemma4}

\begin{lemma}
For any synthesis derivation tree $T_2$, there exists a type derivation tree $T_1$, $T_1 \cong T_2$.
\end{lemma}
\begin{proof}
  Following the translation from synthesis derivation trees to the type derivation trees detailed in \autoref{sec:isomorphism}, the root and child order preservation is easily satisfied. Then we can prove the label correspondence by induction on the size of the synthesis derivation tree $T_2$ (suppose $T_2$'s top-level synthesis rule has the form in \autoref{S-Rule}).

  \begin{itemize}
     \item \textbf{Base case:} The base case occurs when the synthesis derivation tree $T_2$ has only one node, whose synthesis rule has no premises. By the described translation, the node in the type derivation tree $T_1$ is annotated with the typing judgment $P(\theta(\overline{x}))$, if $\theta$ is the synthesized assignment. In this case, the label correspondence is clearly satisfied.
     \item \textbf{Inductive case:} The inductive case arises when the synthesis rule has one or more subgoals. 
     By the induction hypothesis, for the $i$-th subgoal $P_i(\sigma_{i-1}\ldots\sigma_0\sigma(\overline{x}_i)) \leadsto \sigma_i$, we have constructed a type derivation subtree whose root node is labeled with the typing judgment $P_i(\overline{x}_{ic})$, where $\overline{x}_{ic}=\sigma_{i}\ldots\sigma_0\sigma(\overline{x}_i)$. 
     
     Let $\theta_c=\sigma_{n}\ldots\sigma_0\sigma|_{\mvhzc(\overline{x}_0\cup\ldots\cup\overline{x}_n)}$, then we have $\theta_c(\overline{x}_{i})=\sigma_{i}\ldots\sigma_0\sigma(\overline{x}_i)=\overline{x}_{ic}$ for all $i\geq 1$. 
     For the conclusion part, $\theta_c(\overline{x}_{0})=\sigma_0\sigma(\overline{x}_0)=\sigma_0\sigma(\overline{x})=\sigma_n \ldots \sigma_0\sigma(\overline{x})=\theta(\overline{x})$. 
     Thus, under the assignment $\theta_c$, the corresponding typing rule, along with the translated type derivation subtrees, forms a valid type derivation tree $T_1$ whose root node is labeled with $P(\theta(\overline{x}))$.
 \end{itemize}
\end{proof}

\clearpage
\section{Additional Information about Benchmarks}
\label{appendix-benchmark}

\subsection{Java Benchmark}
We have carefully designed a subset of the Java programming language that encompasses the fundamental syntactic constructs, including classes, methods, variables, arrays, generic types, operators, and control flow statements. While excluding advanced Java features such as multithreading, lambda expressions, and complex reflection mechanisms.

The Java subset retains sufficient expressiveness to represent core imperative programming paradigms. Our evaluation shows that it covers approximately 78\% of the Java programs in the original MBJP dataset, which is the Java version of the MBPP dataset and available at \href{https://github.com/amazon-science/mxeval}{amazon-science/mxeval}.
The subset admits tasks whose constructs our toolchain supports; this admission filter is a potential selection bias, mitigated by the two additional benchmarks drawn from independent sources.
The complete grammar specification is presented below:

\begin{tcolorbox}[
    enhanced,
    breakable,
    colback=blue!5!white,
    colframe=blue!50!black,
    title=Java Subset Language Grammar,
    fonttitle=\bfseries,
    boxrule=1pt,
    arc=3pt
]
\begin{Verbatim}[fontsize=\small]
<program>     := <program> <program>
              |  <modifier> <type> <identifier> ;
              |  <modifier> <type> <identifier> = <term> ;
              |  <modifier> <type> <identifier> ( <stmt> ) { <stmt> }
              |  <identifier> ( <stmt> ) { <stmt> }
              |  <stmt>
              |  class <identifier> { <program> }
<stmt>        := empty
              |  <type> <identifier> = <term>;
              |  <type> <identifier>;
              |  if ( <term> ) { <stmt> }
              |  if ( <term> ) { <stmt> } else { <stmt> }
              |  while ( <term> ) { <stmt> }
              |  do { <stmt> } while ( <term> )
              |  for ( <stmt> ; <term> ; <term> ) { <stmt> }
              |  for ( <type> <identifier> : <term> ) { <stmt> }
              |  return <term>;
              |  continue;
              |  break;
              |  <term>;
              |  switch ( <term> ) { <stmt> }
              |  case <term> : <stmt>
              |  <stmt> <stmt>
<term>        := <identifier>
              |  <integer>
              |  <float>
              |  <character>
              |  <string>
              |  true
              |  false
              |  null
              |  <term> = <term>
              |  ( <type> ) <term>
              |  <term> instanceof <type>
              |  <term> ? <term> : <term>
              |  <term> . <identifier>
              |  <term> [ <term> ]
              |  new <type> ( <term> )
              |  new <type>
              |  new <term> [ <term> ]
              |  <term> . <identifier> ( <term> )
              |  <identifier> ( <term> )
              |  <type>
              |  [ <term_list> ]
              |  { <term_list> }
              |  ( <term> )
              |  <term> + <term>
              |  <term> - <term>
              |  <term> * <term>
              |  <term> / <term>
              |  <term> % <term>
              |  - <term>
              |  <term> << <term>
              |  <term> >> <term>
              |  <term> & <term>
              |  <term> | <term>
              |  <term> ^ <term>
              |  ~ <term>
              |  <term> ++
              |  <term> --
              |  ++ <term>
              |  -- <term>
              |  <term> == <term>
              |  <term> != <term>
              |  <term> < <term>
              |  <term> <= <term>
              |  <term> > <term>
              |  <term> >= <term>
              |  <term> && <term>
              |  <term> || <term>
              |  ! <term>
<type>        := int
              |  float
              |  boolean
              |  char
              |  String
              |  null
              |  <type> []
              |  <identifier> <>
              |  <identifier> < <type> >
              |  <identifier> < <type> , <type> >
              |  <type> -> <type>
              |  <identifier>
              |  void
<term_list>   := <term> , <term_list>
              |  <term>
              |  empty
<modifier>    := public | private | protected | static | final | abstract
<identifier>  := [a-zA-Z_][a-zA-Z0-9_]*
<integer>     := [0-9]+
<float>       := [0-9]+.[0-9]+
<character>   := '[^']'
<string>      := ".*"
\end{Verbatim}
\end{tcolorbox}

\subsection{HumanEval-Java and TransCoder-GFG}

\textbf{HumanEval-Java} is the Java port of HumanEval~\cite{DBLP:journals/corr/abs-2107-03374}: short Python-style algorithm snippets (string/list manipulation, digit arithmetic, parity and divisibility tests, a few short geometry/number-theory problems) transcribed into a single Java method per task. \textbf{TransCoder-GFG} reuses the GeeksforGeeks Java programs released with TransCoder~\cite{DBLP:conf/nips/LachauxROC20}: one- or two-sentence specifications (numeric, string, or array algorithms; standard-library types only) paired with a self-contained Java method. \autoref{tab:java-benchmark-stats} summarizes the three Java evaluation splits used in this paper.

\begin{table}[ht]
\centering
\small
\caption{Test-split statistics of the three Java benchmarks.}
\label{tab:java-benchmark-stats}
\begin{tabular}{l r r}
\toprule
Benchmark (train / test) & prompt (words) & gold LoC \\
\midrule
MBJP (608 / 67)~\cite{DBLP:conf/iclr/AthiwaratkunGWL23}             & 63\,(46--137) & 27\,(19--43) \\
HumanEval-Java (146 / 16)~\cite{DBLP:journals/corr/abs-2107-03374}  & 88\,(53--135) & 12\,(5--53)  \\
TransCoder-GFG (414 / 103)~\cite{DBLP:conf/nips/LachauxROC20}       & 65\,(55--105) & 12\,(3--31)  \\
\bottomrule
\end{tabular}
\\[2pt]
\footnotesize Cell entries are \emph{median (range)}; prompt length counts whitespace-separated words on the full prompt (imports + class skeleton + javadoc); gold LoC counts non-blank lines of the canonical reference solution; difficulty labels for TransCoder-GFG (easy / medium / hard) are 53 / 49 / 1.
\end{table}

HumanEval-Java and TransCoder-GFG problems are short single-method snippets (median 12 LoC, 21--35-word prompts), whereas MBJP is the densest and most example-rich benchmark (median 27 LoC, prompt with three executable example cases and a gold solution about twice as long as the additional benchmarks).

\paragraph{Strengthened HumanEval-Java and TransCoder-GFG baselines.}
Initial fine-tuning of the HumanEval-Java and TransCoder-GFG baselines used a learning rate of $5\times10^{-5}$ with last-checkpoint selection; a post-hoc recipe audit found this setting overfits the small training splits and underestimates the base model. The reported baselines instead fine-tune T5Gemma2-2B with learning rate $10^{-5}$ for 30 passes---on the union of the three Java training splits (1{,}168 tasks) for HumanEval-Java, and on the TransCoder-GFG training split plus the MBJP training tasks (1{,}022 tasks) for TransCoder-GFG---selecting the checkpoint on validation holdouts drawn from the respective training splits (32 and 60 tasks, never in any training data). Each holdout-selected configuration was then evaluated once on the fixed test split. The same audit found no benefit from excluding cross-benchmark data, and the strengthened baselines retain the standard beam-10 decoding protocol without few-shot prompting. Per-candidate results and the full recipe-search record are part of the artifact package.

\subsection{SuFu Benchmark}

The SuFu dataset is sourced from the GitHub repository accompanying the SuperFusion research paper (\href{https://github.com/jiry17/SuFu}{jiry17/SuFu}). SuperFusion is an advanced program optimization tool that processes SuFu programs as input and applies a series of transformations to improve computational efficiency. 
The SuFu language features a type system that explicitly specifies those intermediate data structures to be optimized, and all code fragments operating on these data structures.
The SuFu programs in this dataset are drawn from two primary lines of prior work. The first source consists of fusion tasks primarily collected from classical textbooks such as Introduction to Algorithms, while the second originates from research on program restructuring, where the goal is to transform a reference program into specific target forms such as divide-and-conquer or single-pass. 
Each program is supplemented with a natural language description. These descriptions were first generated by the GPT-o3-mini model and then carefully reviewed and edited by human experts to correct inaccuracies and improve clarity.

A key issue with the SuFu language is that when an algorithm description is provided, the initial SuFu program still needs to be manually constructed.
However, manually constructing SuFu programs from algorithmic descriptions is highly challenging and requires considerable domain expertise in functional programming. 
Directly prompting language models to generate SuFu code from natural language descriptions often results in programs that are syntactically or semantically flawed, typically containing multiple type errors. 

The following is an example of a generated SuFu program with type errors:
\begin{tcolorbox}[
    colback=gray!5,
    colframe=black,
    boxrule=0.5pt,
    arc=2pt,
    left=8pt,
    right=8pt,
    top=8pt,
    bottom=8pt,
    title=Natural Language Prompt
]
\textit{Write a function for the mts algorithm, which finds the maximum sum among all possible tail lists of the input list.}
\end{tcolorbox}
\begin{tcolorbox}[
    colback=gray!5,
    colframe=black,
    boxrule=0.5pt,
    arc=2pt
]
\begin{Verbatim}[fontsize=\small, xleftmargin=1em]
tails :: List -> Packed NList
tails Nil = let ts = NCons(Nil, NNil) in
            rewrite (label ts)
tails Cons(_, t)@xs = let ts = tails t in
                      rewrite (label NCons(xs, ts))
                      // there is a type error, `ts` should be `unlabel ts`
mts_label xs = let ts = tails xs in
                rewrite (maximum (map sum (unlabel ts)))
\end{Verbatim}
\end{tcolorbox}

The complete SuFu grammar specification is as follows:
\begin{tcolorbox}[
    colback=blue!5!white,
    colframe=blue!50!black,
    title=SuFu Language Grammar,
    fonttitle=\bfseries,
    boxrule=1pt,
    arc=3pt,
    breakable
]
\begin{Verbatim}[fontsize=\small]
<program>     := <program> <program>
              |  <command>
<command>     := Inductive <Name> = <name> <type> 
                                    | ... 
                                    | <name> <type>;
              |  <Name> = <type>;
              |  <name> = <term>;
<type>        := {<type>, ..., <type>}
              |  <type> -> <type>
              |  Reframe <type>
              |  Unit 
              |  Int 
              |  Bool 
              |  <Name>
<term>        := unit 
              |  <integer> 
              |  true 
              |  false 
              |  <name>
              |  + | - | * | / 
              |  < | <= | > | >= | == 
              |  and | or | not
              |  <term> <term>
              |  \<name>:<type>. <term>
              |  fix <term>
              |  let <name> = <term> in <term>
              |  if <term> then <term> else <term>
              |  {<term>, ..., <term>}
              |  <term>.<integer>
              |  match <term> with <pattern> -> <term> 
                                   | ... 
                                   | <pattern> -> <term> 
                 end
              |  rewrite <term> 
              |  label <term> 
              |  unlabel <term>
<pattern>     := _ 
              |  <name>
              |  <name> <pattern>
              |  {<pattern>, ..., <pattern>}
<Name>        := [A-Z][a-zA-Z0-9_]*
<name>        := [a-z][a-zA-Z0-9_]*
<integer>     := [0-9]+
\end{Verbatim}
\end{tcolorbox}

\clearpage
\section{Failure Analysis}\label{appendix-failure}

This section classifies the failing tasks of the Java evaluation and presents representative cases.
For each benchmark and model, every test task is assigned to one of four mutually exclusive categories:
\textbf{Solved@1} (the top-1 candidate passes all tests);
\textbf{Ranking failure} (some candidate in the top ten passes, but not the top-1 one);
\textbf{All-invalid} (no candidate passes and none of the ten candidates compiles); and
\textbf{Well-typed but wrong} (no candidate passes, but at least one candidate compiles---the failure is semantic, not syntactic).

\begin{table}[ht]
\centering
\begin{threeparttable}
\caption{Task-Level Failure Taxonomy on the Three Java Test Sets}
\label{tab:failure-taxonomy}
\begin{tabular}{
    l
    l
    S[table-format=2.0]
    S[table-format=2.0]
    S[table-format=2.0]
    S[table-format=2.0]
    S[table-format=3.0]
}
\toprule
Benchmark & Model & {Solved@1} & {Ranking} & {All-invalid} & {\shortstack{Well-typed\\wrong}} & {Total} \\
\midrule
\multirow{2}{*}{MBJP (67)}
        & T5Gemma2-2B & 9 & 13 & 4 & 41 & 67 \\
        & \mainname-2B & 17 & 12 & 0 & 38 & 67 \\
\midrule
\multirow{2}{*}{HumanEval-Java (16)}
        & T5Gemma2-2B & 5 & 1 & 0 & 10 & 16 \\
        & \mainname-2B & 8 & 1 & 0 & 7 & 16 \\
\midrule
\multirow{2}{*}{TransCoder-GFG (103)}
        & T5Gemma2-2B & 20 & 14 & 2 & 67 & 103 \\
        & \mainname-2B & 31 & 17 & 2 & 53 & 103 \\
\bottomrule
\end{tabular}
\end{threeparttable}
\end{table}

Counts in \autoref{tab:failure-taxonomy} are derived from per-candidate statuses of the frozen runs; timeouts and unfilled slots count as invalid, and the two TransCoder-GFG all-invalid \mainname-2B tasks include the documented budget-exhaustion (fail-closed) task.
Invalid generation essentially disappears under \mainname{}: the all-invalid category drops from 4/0/2 to 0/0/2 tasks, the residual GFG cases being budget exhaustion rather than ill-formed output.
Ranking failures persist (12--17 tasks on MBJP and TransCoder-GFG): the constrained search often finds a correct derivation but does not rank it first, a property of beam scoring rather than the representation.
The dominant residual mode is well-typed but semantically wrong code: 38/50 unsolved MBJP, 7/8 unsolved HumanEval-Java, and 53/72 unsolved GFG tasks have at least one compilable candidate with incorrect behaviour; this is exactly the boundary of what type-level constraints can guarantee.

\subsection{Representative Case: Invalid Generation (T5Gemma2-2B)}

MBJP task~5 (\emph{tupleStrInt}, converting a string tuple to an integer list) is unsolved by the ordinary baseline with all ten candidates failing to compile.
The top-1 candidate omits the opening brace of the loop body:
\begin{tcolorbox}
\begin{Verbatim}[fontsize=\small]
public static List<Integer> tupleStrInt(String testStr) {
    int count = 0;
    List<Integer> ret = new ArrayList<>();
    for (String s : testStr.split("[, ]+"))
        count = count << 1;
        ret.add(Integer.valueOf(s.trim()));
    }
    return ret;
\end{Verbatim}
\end{tcolorbox}
Without the brace, the loop body is a single statement, the closing brace terminates the method early, and \texttt{javac} reports \texttt{illegal start of type} at the \texttt{return} statement.
The model produces plausible token-level continuations, but nothing prevents such structural corruption.

\subsection{Representative Case: Well-Typed but Wrong (\mainname-2B)}

MBJP task~0 (\emph{oddPosition}, checking that every odd index holds an odd number) is unsolved by \mainname-2B, yet its top-1 candidate compiles and type-checks:
\begin{tcolorbox}
\begin{Verbatim}[fontsize=\small]
public static boolean oddPosition(List<Integer> nums) {
    boolean res = false;
    for (int i = 0; i < nums.size(); i++) {
        if (i % 2 == 0) {
            res = false;
        }
    }
    return res;
}
\end{Verbatim}
\end{tcolorbox}
The candidate has the right shape---a loop over indices with a modulo test---but the body never inspects \texttt{nums.get(i)} and never assigns \texttt{true}; the method returns \texttt{false} for every input, while the tests expect \texttt{true} for two of the three cases.
No typing rule is violated, so no amount of type-level pruning can reject this candidate: the failure is purely semantic.

\subsection{SuFu}

On SuFu, \mainname eliminates compilation errors entirely (CER $=$ 0 in \autoref{tab:model-results}), so every unsolved task falls into the well-typed-but-wrong or ranking category; the fixed-budget and fail-closed behavior is described in the RQ2 discussion.

\clearpage
\section{Statistical Significance Analysis}
\label{appendix-statistics}

This section complements the point estimates in \autoref{tab:model-results} with confidence intervals and paired significance tests. It addresses two questions: whether \mainname's improvements over its base models are statistically distinguishable, and how much uncertainty the small test sets leave around each reported rate.

\subsection{Paired Significance Tests at the 220M Scale}
\label{appendix-statistics-220m}

We pair every CodeT5-220M baseline run with the matching \mainname-220M run on the same tasks for both benchmarks, and compute exact two-sided McNemar tests on pass@1 / pass@10 and exact two-sided sign tests on FSP and on per-task compilation-error fractions (CER).
Both pairs use the same frozen 220M checkpoint, the same beam size, and the same scoring script, so the two rows of \autoref{tab:paired-statistics-220m} are directly comparable across languages.
The corresponding 2B aggregate rates remain in \autoref{tab:model-results}; \autoref{appendix-statistics-2b} reports the paired tests for the two 2B benchmarks whose per-task arrays are preserved.

\begin{table}[ht]
\centering
\begin{threeparttable}
\caption{Paired Significance Tests at the 220M Scale}
\label{tab:paired-statistics-220m}
\begin{tabular}{l l l l l l}
\toprule
Benchmark & Metric & Baseline & \mainname & 95\% interval (baseline / \mainname) & Paired $p$-value \\
\midrule
\multirow{4}{*}{Java}
 & pass@1  & 7/67 (10.45\%)  & 17/67 (25.37\%)  & [5.15,20.03] / [16.49,36.93]   & $2.13{\times}10^{-2}$ \\
 & pass@10 & 14/67 (20.90\%) & 29/67 (43.28\%) & [12.88,32.07] / [32.10,55.19] & $1.50{\times}10^{-3}$ \\
 & FSP     & 8.19            & 6.36            & --                              & $1.25{\times}10^{-2}$ \\
 & CER     & 238/670 (35.52\%) & 3/670 (0.45\%)  & --                              & $1.78{\times}10^{-13}$ \\
\midrule
\multirow{4}{*}{SuFu}
 & pass@1  & 14/58 (24.14\%)  & 22/58 (37.93\%)  & [14.96,36.53] / [26.56,50.80] & $3.86{\times}10^{-2}$ \\
 & pass@10 & 19/58 (32.76\%)  & 27/58 (46.55\%)  & [22.08,45.58] / [34.33,59.20] & $9.23{\times}10^{-2}$ \\
 & FSP     & 7.03             & 5.48             & --                              & $1.58{\times}10^{-2}$ \\
 & CER     & 482/580 (83.10\%) & 0/580 (0.00\%)   & [79.84,85.93] / [0.00,0.66]    & $6.94{\times}10^{-18}$ \\
\bottomrule
\end{tabular}
\begin{tablenotes}
\small
\item Java: MBJP, 67 tasks; SuFu: 58 tasks; ten candidates per task for both.
\item pass@1 and pass@10 use Wilson intervals over tasks; their paired tests are exact two-sided McNemar on the discordant pairs.
\item FSP is the mean first-success rank with unsolved tasks assigned rank 10; its interval is omitted because each benchmark here has only 58 or 67 tasks, which is too few for a normal approximation.
\item CER uses Wilson intervals over returned candidates; the paired test is an exact two-sided sign test on per-task compilation-error fractions.
\item In the Java block all four paired differences favor \mainname{} at the 0.05 significance level; in the SuFu block pass@1, FSP, and CER do, while pass@10 is borderline ($p=9.23{\times}10^{-2}$). CER differences reach the resolution limit of the test.
\end{tablenotes}
\end{threeparttable}
\end{table}

\subsection{Paired Significance Tests at the 2B Scale}
\label{appendix-statistics-2b}

The HumanEval-Java and TransCoder-GFG baselines and the matched \mainname-2B runs both retain per-task records, so \autoref{tab:paired-statistics-2b} repeats the same tests at the 2B scale.
On TransCoder-GFG (103 tasks) the pass@10 improvement is significant ($p=0.013$), as is the candidate-level compilation-error difference ($p<10^{-5}$); pass@1 and FSP fall just outside the 0.05 level ($p=0.061$ and $p=0.096$).
On HumanEval-Java (16 tasks) no paired test reaches significance, which is expected at this test size; the point estimates and the CER gap favor \mainname{} in the same direction, and we treat the 16-task comparison as directional evidence only.
The MBJP and SuFu 2B rows retain aggregate counts only, so no paired test is reported for them.

\begin{table}[ht]
\centering
\begin{threeparttable}
\caption{Paired Significance Tests at the 2B Scale}
\label{tab:paired-statistics-2b}
\begin{tabular}{l l l l l l}
\toprule
Benchmark & Metric & Baseline & \mainname-2B & 95\% interval (baseline / \mainname) & Paired $p$-value \\
\midrule
\multirow{4}{*}{HumanEval-Java (16)}
 & pass@1  & 5/16 (31.25\%)  & 8/16 (50.00\%)  & [14.16,55.60] / [28.00,72.00] & $3.75{\times}10^{-1}$ \\
 & pass@10 & 6/16 (37.50\%)  & 9/16 (56.25\%)  & [18.48,61.36] / [33.18,76.90] & $2.50{\times}10^{-1}$ \\
 & FSP     & 6.50            & 4.75            & --                            & $3.75{\times}10^{-1}$ \\
 & CER     & 24.38\%         & 1.25\%          & --                            & $9.77{\times}10^{-4}$ \\
\midrule
\multirow{4}{*}{TransCoder-GFG (103)}
 & pass@1  & 20/103 (19.42\%) & 31/103 (30.10\%) & [12.94,28.10] / [22.09,39.54] & $6.14{\times}10^{-2}$ \\
 & pass@10 & 34/103 (33.01\%) & 48/103 (46.60\%) & [24.68,42.56] / [37.26,56.18] & $1.25{\times}10^{-2}$ \\
 & FSP     & 7.10             & 5.81             & --                            & $9.61{\times}10^{-2}$ \\
 & CER     & 13.69\%          & 2.72\%           & --                            & $3.39{\times}10^{-6}$ \\
\bottomrule
\end{tabular}
\begin{tablenotes}
\small
\item Ten candidates per task; the same scorers and decoding protocols as \autoref{tab:model-results}.
\item pass@1 and pass@10 use Wilson intervals over tasks and exact two-sided McNemar tests on the discordant pairs.
\item FSP is the mean first-success rank with unsolved tasks assigned rank 10; the interval is omitted because the mean is not used for significance here, and the paired test is an exact two-sided sign test.
\item CER entries are per-task mean compilation-error fractions over the ten nominal slots per task; their paired test is an exact two-sided sign test on the per-task fractions. The HumanEval-Java \mainname-2B value (1.25\%) therefore counts the six ungenerated slots as valid, whereas \autoref{tab:model-results} uses the 154 returned candidates as denominator (1.30\%).
\end{tablenotes}
\end{threeparttable}
\end{table}
